\section{What the bound certifies, and what it does not}
\label{sec:bnd:cert}

Theorem~A (Theorem~\ref{thm:bnd2:reg}) and Theorem~B (Theorem~\ref{thm:bnd2:sub}) bound a
truncation. They do
not bound a rounding error, and no bound in this document does. This section separates the
two experimentally, on the hardest sample available, and states the one certificate that
holds with no hypothesis at all.

\subsection{Sixty random cases, and the classification of every violation}
\label{sec:bnd:cert:60}

The sample is the adversarial audit's: sixty cases drawn with a fixed seed, the shape
uniform over twelve shapes of aspect ratio $1$ to $1020$ plus $(1/32)^3$ and $(1/4)^3$
(thirteen of the fourteen were drawn), the offset uniform in $[-20,20]^3$ with
$\max|\mathbf{D}| \ge 2$, and the frequency uniform over
$\{1,\ 1+0.1\ii,\ 0.37,\ 2+0.2\ii,\ 1+\ii\}$. The reference is an independent $220$-bit
graded Gauss--Legendre volume rule --- no addition theorem, no moments, no face pairs, no
shared code with the library --- whose own convergence over these shapes is a median
$6.0\times10^{-50}$. \emph{bound} is the library's own tail at the $L$ it actually selected,
\emph{actual} is $\max_{ab}|G - G_{\mathrm{ref}}|$, \emph{flr} is
$\varepsilon\max_{ab}|G_{\mathrm{ref}}|$, and \emph{amp} is
$\max_{ab}\sum_l|\text{term}|/|\text{sum}|$ of the $l$-sum the library performs.

\begin{center}\small
\begin{tabular}{@{}lrrrr@{}}
\toprule
over the 60 cases & min & median & max & violations ($<1$)\\
\midrule
bound / actual, shipped library & $0.118$ & $4.02$ & $54.5$ & 15\\
bound / actual, the audited version & $9.96\times10^{-10}$ & $7.48$ & $419$ & 19\\
\midrule
actual $/\,\varepsilon\max|G_{\mathrm{ref}}|$ & $0.636$ & $2.3$ & $\mathbf{16.4}$ & \\
$\max_{ab}|G-G_{\mathrm{ref}}|/\max|G_{\mathrm{ref}}|$ & & & $3.65\times10^{-15}$ & \\
$l$-sum amplification $\sum|\text{term}|/|\text{sum}|$ & $1.27$ & $3.7$ &
$3.14\times10^{3}$ & \\
\bottomrule
\end{tabular}
\end{center}

\noindent
\textbf{All fifteen remaining violations are the rounding floor. None is a truncation
violation.} In every one of them the Theorem~A bound lies \emph{below} the \code{Float64}
evaluation floor of the answer --- bound $1.8\times10^{-22}$ to $1.8\times10^{-19}$ against
$\varepsilon\max|G_{\mathrm{ref}}| = 2\times10^{-22}$ to $3\times10^{-19}$ --- so what it
fails to cover is rounding, not truncation. Their errors run from $0.964$ to $15.6$ times
$\varepsilon\max|G_{\mathrm{ref}}|$ and their $l$-sum amplifications from $1.78$ to $388$:

\begin{center}\footnotesize
\begin{tabular}{@{}llrrrrr@{}}
\toprule
shape & $\mathbf{D}$ & $f$ & $L$ & bnd/act & act/flr & amp\\
\midrule
$r_3$ & $(-14,13,-5)$ & $1$ & 28 & $0.717$ & $15.6$ & $388$\\
$r_3$ & $(12,9,-10)$ & $1+\ii$ & 32 & $0.122$ & $10.5$ & $45.1$\\
$r_2$ & $(17,19,15)$ & $1+\ii$ & 24 & $0.840$ & $6.09$ & $36.2$\\
$\mathrm{pl}_8$ & $(-17,-11,-7)$ & $1$ & 16 & $0.293$ & $5.56$ & $3.50$\\
$\mathrm{sl}_X$ & $(-2,13,1)$ & $1+\ii$ & 16 & $0.635$ & $4.53$ & $2.79$\\
$\mathrm{pl}_8$ & $(20,14,-8)$ & $0.37$ & 14 & $0.118$ & $2.82$ & $2.70$\\
$r_1$ & $(7,-19,10)$ & $2+0.2\ii$ & 16 & $0.119$ & $2.76$ & $14.5$\\
\multicolumn{7}{@{}l@{}}{\emph{(eight further rows at act/flr $\le 2.8$, amp $\le 6.9$)}}\\
\bottomrule
\end{tabular}
\end{center}

\noindent
For comparison, the same sixty cases on the version this one replaces had nine
\emph{genuine} failures --- the answer wrong, not the bound loose --- with \emph{bound /
actual} down to $9.96\times10^{-10}$ and errors up to $1.38\times10^{-4}$ in max-norm, and
four more where the answer was right to $10^{-13}$ but the rounding error reached
$396\,\varepsilon$. Both classes are now empty: $r_3$ at $(18,-14,-10)$ went from
$2.43\times10^{-3}$ to $54.5$, $r_3$ at $(-14,13,-5)$ from $9.60\times10^{-3}$ to $0.717$,
the $\lambda/128$ cube at $(-18,7,-13)$ from $3.86\times10^{-2}$ to $11.5$.

\subsection{The statement this document carries}

\begin{quote}
\emph{The truncation error is certified by the bound: over the sixty cases no violation is
a truncation violation, and the largest truncation-side ratio is
$\mathrm{bound}/\mathrm{actual} = 54.5$. The rounding error is measured, not bounded, at no
more than $X\varepsilon$ times the largest tensor entry with
\[
X = 16.4 ,
\]
attained on an aspect-$87.7$ cell at $f = 2+0.2\ii$ whose $l$-sum amplification is
$3.14\times10^{3}$.}
\end{quote}

\noindent
The corresponding a posteriori rule is $\mathrm{bound} + 20\,\varepsilon\max_{ab}|T_{ab}|$,
which no row of the sixty exceeds. (The audited version needed
$\mathrm{bound}+\text{several hundred }\varepsilon$ and still failed on thirteen rows.)

\subsection{The certificate that needs no hypothesis}
\label{sec:bnd:cert:post}

Everything above compares the bound against $\mathrm{tol}\cdot\mathrm{est}(\mathbf{R})$, and
$\mathrm{est}$ is an a priori scale, not a bound on $\max_{ab}|T_{ab}|$
(Section~\ref{sec:unc}). After the tensor has been evaluated that hypothesis can be
discharged, because both the bound and the answer are in hand:
\begin{equation}
\frac{\max_{ab}|T_{ab}-T^{(L)}_{ab}|}{\max_{ab}|T_{ab}|}
\;\le\; \frac{\mathrm{bnd}}{\max_{ab}|T^{(L)}_{ab}| - \mathrm{bnd}} ,
\label{eq:cert:post}
\end{equation}
which follows from the triangle inequality alone and costs one already-computed number.
Over the $427$ cached references \eqref{eq:cert:post} evaluates to at most
$4.5\times10^{-14}$. \textbf{That is the honest guarantee: a certified max-norm relative
truncation error of $4.5\times10^{-14}$, against a measured one of
$8.6\times10^{-15}$}, with the rounding allowance of the previous subsection on top.

A second, proven lower scale exists and is one keyword away. Away from the source
$\nabla^2 g = -k^2g$, so $\operatorname{tr}T = 2k^2 S$ with
$S = (1/V_{\mathrm{t}})\int_D w\,g(\mathbf{R}+\boldsymbol\delta)\dd\boldsymbol\delta$, hence
$\max_{ab}|T_{ab}| \ge |\operatorname{tr}T|/3$, and $|S|$ is bounded below by its $l = 0$
shell minus the Theorem-A tail with $V^{ab}_l$ replaced by $W_l$ --- everything already in
the shape table, $O(L)$ per offset, and a function of $|\mathbf{R}|$ alone, so the threshold
table survives. Measured, $\mathrm{est}_{\mathrm{lo}}/\max|G_{\mathrm{ref}}|$ reaches
$0.98$, $0.63$ and $0.71$ on the $\lambda/32$, $\lambda/8$ and slender shapes, against
$\mathrm{est}/\max|G_{\mathrm{ref}}| \in [1.10, 6.34]$ above; it costs $4$ to $13\,\%$ in
terms where it applies. It is \emph{vacuous} on a $\lambda/4$ cell, where the $l \ge 2$ tail
of the scalar series exceeds its $l = 0$ shell (the ratio
$|k|^2r_{\mathrm{d}}^2/(18\,j_0(|k|r_{\mathrm{d}}))$ is $0.0065$, $0.14$ and $0.0009$ at
$\lambda/32$, $\lambda/8$ and the slender cell, and $2.86$ at $\lambda/4$): the trace is
direction-free precisely because it discards the anisotropy, and at small $k\mathbf{R}$ the
anisotropic part is the whole of $T$. So a proven, direction-free, $O(1)$-per-offset scale
exists for cells up to about $\lambda/8$ across the diagonal and does not exist by this
argument at $\lambda/4$. The default remains $\mathrm{est}$, with
\eqref{eq:cert:post} as the universal fallback.
